\begin{helper}[Coordinate Marginal Probability]
Fix $m_{sub} \in \mathbb{N}$, a probability $p_{sub} \in \mathbb{R}$ with $0 \le p_{sub} \le 1$, subsets $P_R, P_B \subseteq \mathrm{Fin}(m_{sub})$, a coordinate $x \in \mathrm{Fin}(2m_{sub})$, and target sets $A_R, A_B \subseteq \mathrm{Fin}(m_{sub})$.
Let $q = \noderef{FixedSetProductModelMassFn}(p_{sub}, P_R, P_B)$.
Then the marginal probability that the coordinate $x$ contains $A_R$ in its first component and $A_B$ in its second component is exactly $p_{sub}^{|A_R| + |A_B|}$ if $A_R \subseteq P_R$ and $A_B \subseteq P_B$, and $0$ otherwise.
\end{helper}
\begin{proof}
Let
\[
S=\mathcal P(\mathrm{Fin}(m_{sub}))\times
  \mathcal P(\mathrm{Fin}(m_{sub})),\qquad
J=\mathrm{Fin}(2m_{sub}),
\]
and write
\[
W(v)=\prod_{i\in J}q(v_i).
\]
By \(\noderef{FixedSetProductModelMass}\), the one-coordinate mass \(q\)
satisfies
\[
\sum_{a\in S}q(a)=1.
\]

First isolate the distinguished coordinate \(x\).  Let
\(J_x=J\setminus\{x\}\).  A full assignment \(v:J\to S\) is equivalent to
the pair consisting of its value \(a=v_x\in S\) and its restriction
\(w:J_x\to S\).  The inverse sends a pair \((a,w)\) to the assignment whose
value at \(x\) is \(a\) and whose value at every \(i\ne x\) is \(w_i\).
Under this bijection the product weight factors as
\[
W(v)=q(a)\prod_{i\in J_x}q(w_i).
\]
Therefore finite Fubini gives
\[
\begin{aligned}
&\sum_{v:J\to S}
{\bf 1}_{A_R\subseteq(v_x)_1,\ A_B\subseteq(v_x)_2}W(v)\\
&\qquad =
\sum_{a\in S}
{\bf 1}_{A_R\subseteq a_1,\ A_B\subseteq a_2}q(a)
\left(
\sum_{w:J_x\to S}\prod_{i\in J_x}q(w_i)
\right).
\end{aligned}
\]

The remaining-coordinate sum in parentheses is \(1\).  More generally, for
every finite coordinate set \(T\subseteq J\), define
\[
R(T)=\sum_{w:T\to S}\prod_{i\in T}q(w_i).
\]
We prove \(R(T)=1\) by finite induction on \(T\).  For \(T=\emptyset\), there
is one empty assignment and the empty product is \(1\).  If
\(T=T'\cup\{j\}\) with \(j\notin T'\), then an assignment on \(T\) is
equivalent to a value \(a\in S\) at \(j\) together with an assignment on
\(T'\).  Hence
\[
R(T)=\sum_{a\in S}q(a)\,R(T')
=\left(\sum_{a\in S}q(a)\right)R(T')=1\cdot1=1.
\]
Taking \(T=J_x\), the product marginal reduces to the one-coordinate sum
\[
\sum_{a\in S}
{\bf 1}_{A_R\subseteq a_1,\ A_B\subseteq a_2}q(a).
\]

If \(A_R\not\subseteq P_R\), then any \(a\) with \(A_R\subseteq a_1\) has
\(a_1\not\subseteq P_R\), so \(q(a)=0\) by
\(\noderef{FixedSetProductModelMassFn}\).  The same argument applies if
\(A_B\not\subseteq P_B\).  Thus the one-coordinate sum is \(0\) in the
off-support case.

Assume now that \(A_R\subseteq P_R\) and \(A_B\subseteq P_B\).  Since
\(\noderef{FixedSetProductModelMassFn}\) is zero off the support
\(R\subseteq P_R,\ B\subseteq P_B\), the one-coordinate sum equals
\[
\left(
\sum_{A_R\subseteq R\subseteq P_R}
  p_{sub}^{|R|}(1-p_{sub})^{|P_R|-|R|}
\right)
\left(
\sum_{A_B\subseteq B\subseteq P_B}
  p_{sub}^{|B|}(1-p_{sub})^{|P_B|-|B|}
\right).
\]
For the first factor, reindex each \(R\) uniquely as
\[
R=A_R\cup T,\qquad T\subseteq P_R\setminus A_R.
\]
The two pieces are disjoint, so \(|R|=|A_R|+|T|\), and
\(|P_R|-|R|=|P_R\setminus A_R|-|T|\).  Thus the first factor is
\[
p_{sub}^{|A_R|}
\sum_{T\subseteq P_R\setminus A_R}
  p_{sub}^{|T|}(1-p_{sub})^{|P_R\setminus A_R|-|T|}
=p_{sub}^{|A_R|}(p_{sub}+1-p_{sub})^{|P_R\setminus A_R|}
=p_{sub}^{|A_R|}.
\]
The second factor is identical, with \(A_B\) and \(P_B\) in place of
\(A_R\) and \(P_R\), and equals \(p_{sub}^{|A_B|}\).  Multiplying the two
factors gives
\[
p_{sub}^{|A_R|}p_{sub}^{|A_B|}
=p_{sub}^{|A_R|+|A_B|},
\]
which proves the stated marginal in the on-support case.
\end{proof}
